\section{Expanded Related Work}
\label{app:related-work-continued}

\paragraph{Continual Learning.}
Model weights are updated during pre-training and post-training, but are then often deployed as static artifacts for months to serve traffic. This can quickly lead to a gap between in-weight knowledge and relevant real-world skills and information. \textit{Continual learning} aims to address this very gap \cite{de2021continual, kudithipudi2022biological, hadsell2020embracing, wang2024comprehensive}. Recently, self-distillation methods have claimed to allow for near seamless continual learning on new tasks while avoiding catastrophic forgetting common in other continual learning approaches \cite{shenfeld2026self}. This line of work claims that the policy being updated is minimally changed when trained using on-policy distillation from a teacher shown privileged information.

\paragraph{On-Policy Distillation.}
Distillation methods aim to impart richer knowledge into a student model
(usually from an already trained teacher) compared to standard supervised
learning methods \cite{hinton2015distilling, gou2021knowledge,
agarwal2024policy}. For language modeling tasks, several variants of
distillation exist. \textit{Sequence distillation} uses natural language outputs
of one model to directly fine-tune a student model. \textit{Knowledge
distillation} goes one step further and trains a student using the
probabilities that a teacher assigns to its generated sequence. \textit{On-policy
distillation} has emerged as an alternative to standard knowledge distillation
wherein the sequence being scored, and subsequently used for training, is
generated \textit{on-policy} by the student model. This approach alleviates
observed issues between train-test distribution shift that can arise in standard
knowledge distillation.

\paragraph{Self-Improving Distillation.}
Generally speaking, knowledge distillation methods, including sequence
distillation and on/off-policy distillation, require a stronger teacher in order
to improve the student's performance. This can induce a heavy computational
overhead for researchers and also raises an important question: can a model be
used to improve itself without the need for a stronger teacher? Several recent
approaches have attempted to answer this question in the affirmative.
On-Policy Self-Distillation (OPSD)
\cite{zhao2026selfdistilledreasoneronpolicyselfdistillation} utilizes a single
model architecture to act as both teacher and student: the teacher policy is
conditioned on privileged information, such as a ground-truth answer to a math
problem, and is used to score on-policy generations of a student without the
privileged information. It is worth noting that recent updates to the OPSD
method specifically disable per-token divergence feedback on deliberation
tokens, which is necessary to stabilize training. Self-Distillation Fine-Tuning
(SDFT) \cite{shenfeld2026self} proposes a similar method but focuses on
integrating new knowledge corpora for continual learning rather than verifiable
math questions.

We focus on self-distillation methods in the previous vein: methods that use
ground truth privileged information to directly condition a teacher's
token-level feedback to train a student. However, other flavors of
self-distillation methods may target different objectives and mechanisms. For
example, CRISP \cite{sang2026policy} focuses on reducing thinking trace lengths
by distilling a teacher's concise reasoning into a student model.
Self-distillation has also been attempted where the teacher, instead of being
given privileged information or different steering prompts, is instead sampled
at different temperatures \cite{zhang2026embarrassingly}. Approaches like
Self-Distillation Policy Optimization (SDPO)
\cite{hubotter2026reinforcement} also use privileged information but augment
the teacher policy with environmental feedback like error messages instead of
gold, ground-truth information like OPSD or SDFT.

\paragraph{Self-Distillation Failure Modes.}
Concurrent work \citep{kim2026doesselfdistillationsometimesdegrade} studies why
self-distillation can sometimes degrade mathematical reasoning, and attributes
the degradation to suppression of epistemic verbalization under richer teacher
conditioning. Their analysis emphasizes context richness and task coverage:
richer conditioning produces shorter, more confident traces with fewer
uncertainty expressions, which can help narrow in-domain settings but hurt OOD
math generalization. Viewed through our framework, epistemic-verbalization
suppression is a visible lexical subset of fork suppression: overt uncertainty
markers expose some high-entropy fork positions, but many branch-relevant
decisions occur at ordinary mathematical, connective, or formatting
continuations. Our work is therefore complementary. We focus on long-rollout
thinking models and ask how privileged token-level feedback changes the
high-entropy decision points that support test-time search. Rather than treating
epistemic markers as the primary object of study, we analyze fork-like positions
in the teacher distribution, token-level signal reversal along fixed student
trajectories, and the resulting loss of long-budget test-time compute gains.

\paragraph{Forking and Exploration in Reasoning Traces.}
Reasoning ability from a token-level perspective has been analyzed in several
works that have found that some tokens, called forking tokens, can have an
outsized influence on the downstream success of a reasoning trace
\cite{bigelow2024forking, lin2024critical, vassoyan2025ignore,
zhang2026embarrassingly}. These critical tokens often occur at high-entropy
positions in a model's generation and control where and how the reasoning
branches into different paths and alternative strategies. Our work
contextualizes failures of self-improvement distillation methods in the setting
of forking suppression.

\paragraph{Privileged Feedback in RL.}
Standard RL setups may face exploration bottlenecks on complex reasoning tasks where correct rollouts can be rare. To overcome this, several methods in RL also aim to incorporate privileged information into the exploration stage. For example, Privileged On-Policy Exploration (POPE) \cite{qu2026pope} and PrefixRL \cite{setlur2026reuse} utilize some privileged information to guide on-policy exploration. Hybrid approaches like HDPO \cite{ding2026hdpo} augment standard RL training with targeted privileged self-distillation in certain cases. Similar frameworks like Self-Distilled RLVR (RLSD) \cite{yang2026self} incorporate self-distillation to guide update magnitudes of standard RLVR training. Our work sets aside RL training mechanisms to focus on the token-level dynamics of privileged distillation---a harm that hybrid approaches may be able to mitigate.
